\section{Proofs and Theoretical Details}\label{app:proofs}

This appendix restates the stylized model M1 in full, derives the closed-form risks of \Cref{sec:theory} step by step, proves Propositions 1--3, gives the formal statement of Proposition~2$'$ (the constrained-least-squares refinement), and documents the parameter mapping that renders the theory curves of \Cref{fig:synth} free of fitted constants.

\subsection{Model M1: formal statement}\label{app:proofs:m1}

A forecasting instance consists of a lookback window of $w$ steps and a target $h$ steps ahead. Only \emph{level restoration} is compared; the predictable part of the shape component is handled identically by all strategies and is absorbed into a common residual (Assumption~A3 of \Cref{sec:theory}). Each instance draws the following random variables, assumed \emph{mutually independent} and Gaussian:
\begin{center}
\small
\begin{tabular}{@{}lll@{}}
\toprule
Symbol & Distribution & Role \\
\midrule
$g_T$ & $\mathcal{N}(0,\lambda V)$ & covariate-explained target-time level (\cnm{} observes it) \\
$v_T$ & $\mathcal{N}(0,(1-\lambda)V)$ & unexplained target-time level \\
$\Delta$ & $\mathcal{N}(0,\sdel)$ & covariate-orthogonal train$\to$test shift \\
$\delta_x$ & $\mathcal{N}(0,\sx)$ & within-horizon covariate ramp; $g_w = g_T - \delta_x$ \\
$\delta_u$ & $\mathcal{N}(0,h\sigma_u^2)$ & within-horizon unexplained drift; $v_w = v_T - \delta_u$ \\
$\bar\varepsilon$ & $\mathcal{N}(0,\sigma_z^2/w)$ & window-mean shape noise \\
$e$ & $\mathcal{N}(0,\sest)$ & \cnm{} first-stage error, $\hat m = g_T + e$ \\
$\zeta$ & $\mathcal{N}(0,\sigma_\varepsilon^2)$ & common shape residual \\
\bottomrule
\end{tabular}
\end{center}
The derived quantities are
\begin{equation}\label{eq:app-derived}
  L = \Delta + g_T + v_T, \qquad
  M = \Delta + g_w + v_w, \qquad
  \bar y = M + \bar\varepsilon, \qquad
  y = L + \zeta ,
\end{equation}
i.e., the target-time level, the window-time level, the observed window mean, and the target. Here $V = \var(g_T + v_T)$ is the total level variance and $\lambda = \var(g_T)/V \in [0,1]$ is the exogenously explained share, matching $R^2_{\mathrm{level}}$ in \Cref{eq:lambda} and expanding the generative decomposition \Cref{eq:model} at the instance level. The shift $\Delta$ is \emph{common} to window and target (the $\Delta$ row of the table above, part of the M1 specification), which is what instance normalization exploits. All restoration rules are fitted on the training segment; without loss of generality the train-set global mean is zero (Assumption~A4).

All MSE computations below rest on one elementary fact. If $X_1,\dots,X_k$ are mutually independent, zero-mean, and square-integrable, then for any signs $s_i \in \{-1,+1\}$,
\begin{equation}\label{eq:app-orth}
  \mathbb{E}\Big[\Big(\textstyle\sum_{i=1}^{k} s_i X_i\Big)^{2}\Big]
  \;=\; \sum_{i=1}^{k} \var(X_i),
\end{equation}
because every cross term satisfies $\mathbb{E}[X_i X_j] = \mathbb{E}[X_i]\,\mathbb{E}[X_j] = 0$ for $i \neq j$. Each estimator's error below is exactly such a signed sum of independent model variables, so each MSE is read off as a sum of variances.

\subsection{Derivation of the closed-form risks}\label{app:proofs:mse}

\paragraph{\rawm{} (global z-score only)}
The restoration rule is $\hat y = 0$ (the train global mean). The error decomposes as
\begin{align}
  y - \hat y \;=\; L + \zeta \;=\; \Delta + g_T + v_T + \zeta ,
\end{align}
a sum of four independent zero-mean terms, so by \Cref{eq:app-orth}
\begin{equation}\label{eq:app-mseraw}
  \mse_{\rawm} \;=\; \sdel + \lambda V + (1-\lambda)V + \sigma_\varepsilon^2
  \;=\; V + \sdel + \sigma_\varepsilon^2 .
\end{equation}

\paragraph{\inm{} (instance normalization)}
The rule restores the observed window mean, $\hat y = \bar y$. Using \Cref{eq:app-derived},
\begin{align}
  y - \bar y
  &= (L + \zeta) - (M + \bar\varepsilon)
   = (g_T - g_w) + (v_T - v_w) - \bar\varepsilon + \zeta \notag\\
  &= \delta_x + \delta_u - \bar\varepsilon + \zeta ,
\end{align}
where the shift $\Delta$ cancels because it enters $L$ and $M$ identically. By \Cref{eq:app-orth},
\begin{equation}\label{eq:app-msein}
  \mse_{\inm} \;=\; \sx + h\sigma_u^2 + \sigma_z^2/w + \sigma_\varepsilon^2 .
\end{equation}
Neither $V$ nor $\sdel$ appears: removing and restoring the window level makes \inm{} immune to level variance and distribution shift, at the price of the within-horizon displacement $\sx + h\sigma_u^2$ and the window-noise floor $\sigma_z^2/w$.

\paragraph{\cnm{} (conditional normalization)}
The oracle rule restores the explained target-time level, $\hat y = g_T$, so
\begin{equation}\label{eq:app-msecnor}
  y - g_T = v_T + \Delta + \zeta
  \qquad\Longrightarrow\qquad
  \mse_{\cnm\text{-oracle}} = (1-\lambda)V + \sdel + \sigma_\varepsilon^2 .
\end{equation}
The estimated rule restores $\hat m = g_T + e$, where the first-stage error $e$ is independent of all other variables (the first stage is fitted on separate data and evaluated out of sample), so
\begin{equation}\label{eq:app-msecn}
  y - \hat m = v_T + \Delta - e + \zeta
  \qquad\Longrightarrow\qquad
  \mse_{\cnm\text{-est}} = (1-\lambda)V + \sdel + \sest + \sigma_\varepsilon^2 .
\end{equation}
No within-horizon ramp appears: the future covariate path pins the level at target time. These four expressions are the rows of the risk table in \Cref{sec:theory}.

\subsection{Proof of Proposition 2 (crossover)}\label{app:proofs:prop2}

\begin{proof}[Proof of \Cref{prop:crossover}]
The common residual $\sigma_\varepsilon^2$ cancels from every comparison.
\emph{(i)} Subtracting \Cref{eq:app-mseraw} from \Cref{eq:app-msein},
\begin{align}
  \mse_{\inm} - \mse_{\rawm}
  \;=\; \sx + h\sigma_u^2 + \sigma_z^2/w \;-\; V - \sdel ,
\end{align}
which is negative exactly when \Cref{eq:inraw} holds.
\emph{(ii)} Subtracting \Cref{eq:app-msein} from \Cref{eq:app-msecn},
\begin{align}
  \mse_{\cnm\text{-est}} - \mse_{\inm}
  \;=\; (1-\lambda)V + \sdel + \sest \;-\; \sx - h\sigma_u^2 - \sigma_z^2/w ,
\end{align}
so $\mse_{\cnm\text{-est}} < \mse_{\inm}$ if and only if
\begin{equation}\label{eq:app-cross}
  (1-\lambda)V \;<\; B, \qquad
  B \;\equiv\; \sx + h\sigma_u^2 + \sigma_z^2/w - \sdel - \sest .
\end{equation}
Dividing by $V > 0$ and rearranging gives $\lambda > 1 - B/V = \lamstar$, which is \Cref{eq:lamstar}. If $B \le 0$ then \Cref{eq:app-cross} fails for every $\lambda \in [0,1]$ (its left side is nonnegative), equivalently $\lamstar \ge 1$ and \inm{} dominates; if $B \ge V$ then \Cref{eq:app-cross} holds for every $\lambda \in [0,1]$ up to the boundary, equivalently $\lamstar \le 0$ and \cnm{}-est dominates.
\end{proof}

\begin{remark}[Coupled parameters and the numerical root]\label{rem:app-numeric}
\Cref{eq:lamstar} presumes that $B$ does not depend on $\lambda$. Under the synthetic mapping of \Cref{tab:app-map}, however, both $\sx = \lambda c_h$ (the ramp scales with the explained share, $c_h \ge 0$) and the drift term $(1-\lambda)h\sigma_{u,\mathrm{dgp}}^2$ are $\lambda$-dependent. The difference $D(\lambda) \equiv \mse_{\cnm\text{-est}}(\lambda) - \mse_{\inm}(\lambda)$ is then affine in $\lambda$ with slope $-\big(V + c_h - h\sigma_{u,\mathrm{dgp}}^2\big)$, strictly negative at all settings used in this paper; more general couplings (e.g., $\sx$ evaluated numerically from a simulated covariate path) need not be affine. The implementation therefore locates $\lamstar$ as the root of $D$ on $[0,1]$ by bisection (tolerance $10^{-10}$) whenever a sign change exists. As a numerical check of the closed form itself: at the baseline parameters of \Cref{fig:dominance} ($V=1$, $w=96$, $\sigma_z=1$, $\sigma_u^2=0.0036$, $\sest = 0.02$, $\sdel = 0$, $\sx = 0$), \Cref{eq:lamstar} evaluates to $\lamstar = 1 - (24 \cdot 0.0036 + 1/96 - 0.02) = 0.923$ at $h=24$, to $0.664$ at $h=96$, and to a negative value at $h=336$ (where $h\sigma_u^2 = 1.2096 > V$), so \cnm{}-est dominates at every $\lambda$.
\end{remark}

\subsection{Proof of Proposition 3 (horizon monotonicity)}\label{app:proofs:prop3}

\begin{proof}[Proof of \Cref{prop:horizon}]
For either \cnm{} variant write $c \in \{0, \sest\}$ for the ($h$-independent) first-stage term. From \Cref{eq:app-msein,eq:app-msecnor,eq:app-msecn},
\begin{equation}\label{eq:app-gap}
  G(h) \;\equiv\; \mse_{\inm}(h) - \mse_{\cnm}(h)
  \;=\; \sx + h\sigma_u^2 + \sigma_z^2/w - (1-\lambda)V - \sdel - c ,
\end{equation}
in which only the first two terms depend on $h$. Treating $h$ as continuous,
\begin{equation}
  \frac{\partial G}{\partial h} \;=\; \sigma_u^2 + \frac{\partial\, \sx}{\partial h} \;>\; 0 ,
\end{equation}
since $\sigma_u^2 > 0$ and $\sx$ is nondecreasing in $h$: the ramp variance aggregates covariate-driven displacement over the horizon (under the AR(1) mapping of \Cref{tab:app-map}, $s_x^2(h) \propto 1 - a^h/\bar\rho$ with $a \in (0,1)$, increasing in $h$). For discrete horizons the same computation gives $G(h+1) - G(h) = \sigma_u^2 + \big[s_x^2(h{+}1) - s_x^2(h)\big] > 0$, which is \Cref{eq:prop3}. Finally, differentiating \Cref{eq:lamstar} in $h$ yields $\partial \lamstar / \partial h = -\big(\sigma_u^2 + \partial \sx/\partial h\big)/V < 0$, so the crossover threshold falls monotonically as the horizon grows.
\end{proof}

\subsection{Proof of Proposition 1 (representational deficiency)}\label{app:proofs:prop1}

\begin{proof}[Proof of \Cref{prop:deficiency}]
Let $(\bar y, g)$ be centered and jointly Gaussian, and let $y = a\bar y + b g + \kappa\,\bar y g + \eta$ with $\eta$ independent of $(\bar y, g)$, $\mathbb{E}\eta = 0$, $\var(\eta) = \sigma_\eta^2$. Write $\mathcal{F} = \{c_0 + c_1 \bar y + c_2 g\}$ for the interaction-free affine class.

\emph{Step 1: the RevIN-reachable class is a strict subclass of $\mathcal{F}$.}
The normalize--restore pipeline as implemented computes, for \emph{any} backbone $f$,
$\hat y = \bar y + s\, f\big((x - \bar y \mathbf{1})/s,\, g\big)$ with $\bar y, s$ the
lookback-window mean and standard deviation \citep{kim2022revin}. RevIN's learnable affine
parameters $(\gamma,\beta)$ act on $f$'s input at the normalize step and are undone on
$f$'s output at the restore step, so they \emph{conjugate} $f$ rather than cancel:
writing $f'(u,g) = \gamma^{-1}\bigl(f(\gamma u + \beta,\, g) - \beta\bigr)$ returns the
composed map to the displayed form with $f'$ in place of $f$, leaving the reachable set
unchanged. Since the sample variance is shift-invariant, $s$ is a function of the centred
window alone, so the arguments of $f$ are exactly invariant to $x \mapsto x + c\mathbf 1$:
the composed map is \emph{additively separable in $\bar y$ with unit slope}, whatever $f$
is and whatever the joint law of $(\bar y, g)$.

We now condition on $s$; this is an added hypothesis, not a consequence of A1--A4, and it
is needed because with $s$ varying across instances the reachable set contains the
directions $s$ and $s\,g$, which $\mathcal{F}$ does not, so the two classes would be
non-nested. At fixed $s$, an affine backbone $f(u,g) = w^\top u + c_2' g + c_0'$ gives
$\hat y = \bar y + w^\top(x - \bar y \mathbf 1) + s c_2' g + s c_0'$, which reparameterizes
to $\mathcal{R}_{\mathrm{full}} = \{c_0 + \bar y + w^\top(x - \bar y \mathbf 1) + c_2 g\}$,
against $\mathcal{F}_{\mathrm{full}} = \{c_0 + c_1\bar y + w^\top(x - \bar y \mathbf 1) + c_2 g\}$
for \rawm{}. The centred-window term $w^\top(x - \bar y \mathbf 1)$ carries zero coefficient
on $\bar y$ and is common to both, so it cancels from every comparison below and we may
work with $\mathcal{R} = \{c_0 + \bar y + c_2 g\} \subsetneq \mathcal{F}$ without loss;
this is why the argument does not require the instance to be summarized by $(\bar y, g)$
and applies to backbones that consume the whole window, as ours do. Since the inclusion
is strict, the lower bound of Steps 2--4 applies to $\mathcal{R}$ \emph{a fortiori};
Step 6 computes the strict gap and Step 7 the unconstrained limit.

\emph{Step 2: reduction to a projection.}
For any $f \in \mathcal{F}$, independence of $\eta$ gives $\mathbb{E}(y - f)^2 = \sigma_\eta^2 + \mathbb{E}\big(a\bar y + bg + \kappa \bar y g - f\big)^2$. Since $a\bar y + bg \in \mathcal{F}$ can be absorbed into $f$,
\begin{equation}
  \min_{f \in \mathcal{F}} \mathbb{E}(y - f)^2 - \sigma_\eta^2
  \;=\; \kappa^2 \min_{c_0,c_1,c_2} \mathbb{E}\big(\bar y g - c_0 - c_1 \bar y - c_2 g\big)^2
  \;=\; \kappa^2\, \mathrm{dist}^2\big(\bar y g,\, \mathcal{F}\big) ,
\end{equation}
while the Bayes risk is exactly $\sigma_\eta^2$; the excess risk is $\kappa^2 \mathrm{dist}^2(\bar y g, \mathcal{F})$ in $L^2$.

\emph{Step 3: the interaction is orthogonal to the linear directions.}
Third moments of a centered Gaussian vector vanish, so
\begin{equation}
  \mathbb{E}\big[(\bar y g)\, \bar y\big] = \mathbb{E}\big[\bar y^2 g\big] = 0,
  \qquad
  \mathbb{E}\big[(\bar y g)\, g\big] = \mathbb{E}\big[\bar y g^2\big] = 0 .
\end{equation}
Moreover $\langle \mathbf{1}, \bar y\rangle = \mathbb{E}\bar y = 0$ and $\langle \mathbf{1}, g\rangle = 0$, so $\mathrm{span}\{\mathbf 1\} \perp \mathrm{span}\{\bar y, g\}$ and the projection of $\bar y g$ onto $\mathcal{F}$ reduces to its projection onto the constants:
$\Pi_{\mathcal{F}}(\bar y g) = \mathbb{E}[\bar y g] = \covop(\bar y, g)$.

\emph{Step 4: Isserlis' theorem.}
For centered jointly Gaussian variables one has $\mathbb{E}[\bar y^2 g^2] = \var(\bar y)\var(g) + 2\covop(\bar y, g)^2$, whence
\begin{equation}\label{eq:app-isserlis}
  \mathrm{dist}^2(\bar y g, \mathcal{F})
  = \mathbb{E}[\bar y^2 g^2] - \covop(\bar y,g)^2
  = \var(\bar y)\,\var(g) + \covop(\bar y, g)^2
  = \var(\bar y g) ,
\end{equation}
which proves \Cref{eq:prop1}.

\emph{Step 5: evaluation under M1.}
From \Cref{eq:app-derived}, $\bar y = \Delta + (g_T - \delta_x) + (v_T - \delta_u) + \bar\varepsilon$, and every summand except $g_T$ is independent of $g_T$; therefore $\covop(\bar y, g_T) = \var(g_T) = \lambda V$, while $\var(g_T) = \lambda V$ and $\var(\bar y) = V + \sx + h\sigma_u^2 + \sigma_z^2/w + \sdel$. The excess risk equals $\kappa^2\big[\lambda V \var(\bar y) + \lambda^2 V^2\big]$, strictly increasing in $\lambda$: the more exogenously driven the level, the larger the irreducible penalty of any predictor confined to affine restoration.

\emph{Step 6: the pinned-slope gap, \Cref{eq:prop1b}.}
Fix $c_1 = 1$ and subtract $\bar y$ from both sides: the residual to be fitted by
$\mathcal{R}$ is $y - \bar y = (a-1)\bar y + b g + \kappa\, \bar y g + \eta$. By Step 3 the
centred interaction $\bar y g - \covop(\bar y, g)$ is orthogonal to
$\mathrm{span}\{\mathbf 1, \bar y, g\}$ and hence to
$\mathrm{span}\{\mathbf 1, g\} \subset \mathrm{span}\{\mathbf 1, \bar y, g\}$, so its
contribution separates and equals $\kappa^2 \var(\bar y g)$ exactly as in Step 4---the same
value $\mathcal{F}$ pays. What remains is
\begin{equation}
  \min_{c_0, c_2} \mathbb{E}\big[(a-1)\bar y + b g - c_0 - c_2 g\big]^2
  \;=\; (a-1)^2 \, \mathrm{dist}^2\big(\bar y,\, \mathrm{span}\{\mathbf 1, g\}\big)
  \;=\; (a-1)^2 \var(\bar y)\big(1 - \varrho^2\big),
\end{equation}
since $b g$ is absorbed by $c_2$ and the residual of $\bar y$ after regression on $g$ has
variance $\var(\bar y)(1-\varrho^2)$. The corresponding minimum over $\mathcal{F}$ is
attained at $c_1 = a$ and leaves no such term. Subtracting gives \Cref{eq:prop1b}. Two
remarks. First, the gap vanishes exactly when $a = 1$---the level persists across the
horizon---which is the assumption instance normalization encodes, so the proposition
quantifies the cost of that assumption being wrong rather than asserting it always is.
Second, it also vanishes as $\varrho^2 \to 1$: a covariate perfectly collinear with the
window mean can stand in for it, so the pinned slope costs nothing. Under M1 the
risk-minimizing $a$ is below one whenever the window mean is a noisy or stale estimate of
the target-time level---precisely the $\sigma_z^2/w$ and $\sx + h\sigma_u^2$ terms of
\Cref{eq:app-msein}---so \Cref{eq:prop1b} is the representational image of the same trade
that \Cref{eq:lamstar} balances.

\emph{Step 7: the unconstrained limit, \Cref{eq:prop1c}.}
Let the backbone range over all measurable functions. \rawm{} then has access to
$\mathbb{E}[y \mid \bar y, g]$ and attains the Bayes risk, so $\Delta\mse_{\rawm} = 0$.
By Step 1 the \inm{} class is $\{\bar y + \varphi(g)\}$ with $\varphi$ arbitrary (the
centred window is independent of the level in M1 and adds nothing), whose minimizer is
$\varphi^\star(g) = \mathbb{E}[\,y - \bar y \mid g\,]$. Writing
$y - \bar y = (a - 1 + \kappa g)\bar y + b g + \eta$ and conditioning on $g$,
\begin{equation}
  \Delta\mse_{\inm}
  = \mathbb{E}\Bigl[(a-1+\kappa g)^2 \var(\bar y \mid g)\Bigr]
  = \Bigl[(a-1+\kappa\,\mathbb{E} g)^2 + \kappa^2\var(g)\Bigr]\var(\bar y)\bigl(1-\varrho^2\bigr),
\end{equation}
using $\var(\bar y \mid g) = \var(\bar y)(1-\varrho^2)$ for jointly Gaussian $(\bar y, g)$
and $\mathbb{E}[(a-1+\kappa g)^2] = (a-1+\kappa \mathbb{E} g)^2 + \kappa^2\var(g)$. With
$g$ centred this is \Cref{eq:prop1c}. Comparing with Step 6, the \inm-versus-\rawm{}
difference rises from $(a-1)^2\var(\bar y)(1-\varrho^2)$ in the affine case to
\Cref{eq:prop1c}: the interaction term, a \emph{shared} cost when both classes are affine,
becomes an \inm-only cost once the backbone is rich enough for \rawm{} to represent it.

\emph{Hypotheses, and where each is used.} Step 4 is the only place joint Gaussianity is
load-bearing: it is Isserlis' fourth-moment identity, and \Cref{eq:prop1} understates the
interaction cost under heavier-tailed elliptical laws with matched first two moments.
Steps 6 and 7 need only vanishing third cross-moments, i.e.\ $\mathbb{E}[\bar y^2 g] =
\mathbb{E}[\bar y g^2] = 0$; centring of $\bar y$ is immaterial to them, but centring of
$g$ is not---the general form of both is $(a - 1 + \kappa\,\mathbb{E} g)^2$ in place of
$(a-1)^2$. The Monte Carlo verification of \Cref{eq:prop1,eq:prop1b,eq:prop1c} over a grid
of $(a, b, \kappa, \varrho, \var(\bar y), \var(g))$, including non-centred and
heavy-tailed cells, is released as \texttt{experiments/g7\_prop1\_verify.py}.
\end{proof}

\subsection{Proposition 2$'$: class-optimal linear predictors}\label{app:proofs:prop2prime}

M1's estimators are pure restoration rules; a trained linear backbone additionally regresses on its input window. Following \citet{toner2024linear}, each normalization defines an affine pre/post transform around the backbone, so the composite is a \emph{constrained} affine predictor whose optimum has a closed form. Fix a lookback $L$, horizon $h$, and an arm-specific restoration series $\ell_t$ ($\ell \equiv 0$ for \rawm{} and RevIN; $\ell = m$ for \cnm{}-oracle; $\ell = \hat m$ for \cnm{}-est). Let $r_t = (y_t - \ell_t - \mu_r)/\sigma_r$ be the residual series standardized with training-segment statistics, let $x \in \mathbb{R}^L$ denote its lookback window with mean $\bar x = \tfrac1L \mathbf 1_L^\top x$, and let the predictor output all $h$ steps directly. The induced classes are
\begin{align}
  \mathcal{F}_{\rawm} &= \big\{\, x \mapsto Wx + b \,\big\}, \notag\\
  \mathcal{F}_{\textsc{RevIN}} &= \big\{\, x \mapsto W(x - \bar x \mathbf 1_L) + \bar x \mathbf 1_h + b \,\big\}, \label{eq:app-revinclass}\\
  \mathcal{F}_{\cnm} &= \big\{\, x \mapsto \sigma_r (Wx + b) + \mu_r + \ell_{t+1:t+h} \,\big\},\notag
\end{align}
with $W \in \mathbb{R}^{h\times L}$, $b \in \mathbb{R}^h$, and predictions evaluated in the global z-score scale of $y$ (train statistics).

\medskip\noindent\textbf{Proposition 2$'$ (formal).}
\emph{(i) Reparameterization. $\mathcal{F}_{\textsc{RevIN}} = \{x \mapsto W'x + b : W'\mathbf 1_L = \mathbf 1_h\}$, the affine class with unit row sums (full pass-through of the window level).
(ii) Class optima. For each class, the training-segment least-squares problem has the closed-form solution $W^{*} = (X^\top X + \rho I)^{-1} X^\top Y$ over the class parameterization (intercept included; ridge $\rho = 10^{-6}$ for numerical conditioning only); its test-segment MSE is the exact class-optimal risk that stochastic-gradient training of a linear backbone approximates. Define $\lamstar_{\mathrm{OLS}}(h,L)$ as the root in $\lambda$ of $\mse^{\mathrm{ols}}_{\textsc{RevIN}}(\lambda) - \mse^{\mathrm{ols}}_{\cnm\text{-est}}(\lambda)$ along the synthetic family of \Cref{sec:synthetic}.
(iii) M1 is an upper bound. Each M1 restoration rule belongs to its class, so $\mse^{\mathrm{ols}}_{\mathrm{arm}}(\lambda) \le \mse^{\mathrm{M1}}_{\mathrm{arm}}(\lambda)$ pointwise for the population risk under the training distribution; and if the class-optimization gain of \cnm{}-est weakly exceeds that of \inm{} at $\lambda = \lamstar_{\mathrm{M1}}$ while the class-optimal difference is decreasing in $\lambda$, then $\lamstar_{\mathrm{OLS}} \le \lamstar_{\mathrm{M1}}$.}

\begin{proof}
\emph{(i)} Expanding \Cref{eq:app-revinclass}, $W(x - \bar x\mathbf 1_L) + \bar x \mathbf 1_h + b = W'x + b$ with $W' = W\big(I_L - \tfrac1L \mathbf 1_L\mathbf 1_L^\top\big) + \tfrac1L \mathbf 1_h \mathbf 1_L^\top$, and $W'\mathbf 1_L = W\mathbf 1_L - W\mathbf 1_L + \mathbf 1_h = \mathbf 1_h$. Conversely, any $W'$ with $W'\mathbf 1_L = \mathbf 1_h$ is realized by choosing $W = W'$: then $W'(I_L - \tfrac1L \mathbf 1_L\mathbf 1_L^\top) + \tfrac1L\mathbf 1_h\mathbf 1_L^\top = W' - \tfrac1L\mathbf 1_h \mathbf 1_L^\top + \tfrac1L\mathbf 1_h\mathbf 1_L^\top = W'$.
\emph{(ii)} Each class is affine in its parameters, so the training risk is a convex quadratic whose minimizer solves the normal equations; the stated ridge form is its numerically conditioned solution. No stochastic optimization enters the theory values.
\emph{(iii)} Membership: $\hat y = 0$ is $(W,b) = (0,0)$ in $\mathcal{F}_{\rawm}$; the window-mean rule $\hat y = \bar x\mathbf 1_h$ is $(W,b) = (0,0)$ in \Cref{eq:app-revinclass}, equivalently $W' = \tfrac1L \mathbf 1_h\mathbf 1_L^\top$ (unit row sums); the \cnm{} rule $\hat y = \ell_{t+1:t+h}$ (level scale) is $(W,b) = (0,0)$ in $\mathcal{F}_{\cnm}$. Class optima therefore weakly improve on the M1 rules, arm by arm. Writing $G_{\mathrm{arm}}(\lambda) = \mse^{\mathrm{M1}}_{\mathrm{arm}} - \mse^{\mathrm{ols}}_{\mathrm{arm}} \ge 0$ and $D_{\mathrm{ols}} = D_{\mathrm{M1}} - (G_{\cnm\text{-est}} - G_{\inm})$ for the \cnm{}--\inm{} differences, the premise $G_{\cnm\text{-est}} \ge G_{\inm}$ at $\lamstar_{\mathrm{M1}}$ gives $D_{\mathrm{ols}}(\lamstar_{\mathrm{M1}}) \le D_{\mathrm{M1}}(\lamstar_{\mathrm{M1}}) = 0$; if $D_{\mathrm{ols}}$ is decreasing, its root satisfies $\lamstar_{\mathrm{OLS}} \le \lamstar_{\mathrm{M1}}$.
\end{proof}

\begin{remark}[Why the gain asymmetry holds]\label{rem:app-gain}
$\mathcal{F}_{\rawm}$ and $\mathcal{F}_{\cnm}$ contain the \emph{level-tracking direction} $W = \tfrac1L \mathbf 1_h \mathbf 1_L^\top$, which reproduces the residual window mean. Because the unexplained level component is serially correlated across window and target ($\delta_u$ is a random-walk increment), least squares places weight along this direction and strictly reduces risk relative to the pure rules $W = 0$ --- implicit, partial instance normalization. The RevIN class, by contrast, already \emph{enforces} unit pass-through of the window level (part (i)), and the M1 rule is near class-optimal within it, so its gain is comparatively small. This is the premise of part (iii), and it is verified empirically in every $(h,L)$ cell of \Cref{sec:synthetic}: $\lamstar_{\mathrm{M1}} \approx 0.27$--$0.28$ at $h=24$, while the class-optimal crossings --- matched by trained RLinear models \citep{li2023rlinear} to within $0.015$ (maximum absolute deviation across the six $(h,L)$ cells) --- lie an order of magnitude lower. Note also that, since \rawm{} and RevIN share the same residual series ($\ell \equiv 0$), part (i) exhibits $\mathcal{F}_{\textsc{RevIN}} \subsetneq \mathcal{F}_{\rawm}$: in distribution, the unconstrained class weakly dominates, and any out-of-sample advantage of RevIN must come from train--test shift --- the function-class expression of its original motivation \citep{kim2022revin}.
\end{remark}

\subsection{Parameter mapping: from theory to synthetic experiment}\label{app:proofs:mapping}

The theory curves in \Cref{fig:synth} contain \emph{no fitted constants}: every M1 parameter is computed from the data-generating process or measured from the first stage, per \Cref{tab:app-map}.

\begin{table}[t]
\centering
\caption{M1 $\leftrightarrow$ synthetic DGP $\leftrightarrow$ real-data correspondence. The backbone is fixed to a single linear layer (RLinear) to match the linear--Gaussian function class; see \Cref{tab:audit} for the full correspondence audit.}
\label{tab:app-map}
\begin{tabular}{@{}lp{5.6cm}p{4.2cm}@{}}
\toprule
M1 quantity & Synthetic DGP (\Cref{sec:synthetic}) & Real-data counterpart \\
\midrule
$\lambda$ & mixing weight in $m_t = \sqrt{\lambda}\,\beta^\top \tilde x_t + \sqrt{1-\lambda}\,u_t$; population $R^2_{\mathrm{level}} = \lambda$ by construction (verified per series) & $\lps$ (\Cref{sec:lps}) \\
$V$ & $\var(m_t)$ at window-mean scale & variance of the window-mean level \\
$h\sigma_u^2$ & $(1-\lambda)\, h\, \sigma_{u,\mathrm{dgp}}^2$ (the DGP scales $u$ by $\sqrt{1-\lambda}$) & unobserved drift \\
$\sx$ & variance of the $h$-step change in the explained level under AR(1) covariates ($a = 0.9$), of the form $2\lambda V\,(1 - a^h/\bar\rho)$; evaluated numerically & weather ramps announced by NWP \\
$\sigma_z^2/w$ & $\sigma_z^2\, \tau / w$ with $\tau$ from \Cref{eq:app-tau} & window-mean noise of the shape component \\
$\sdel$ & accumulated train$\to$test displacement of $u$ & secular level shifts (e.g., demand growth) \\
$\sest$ & measured out-of-sample error of the first-stage LightGBM level model \citep{ke2017lightgbm} & idem, measured before training \\
\bottomrule
\end{tabular}
\end{table}

Two entries deserve derivation. First, the \emph{effective window noise}: the shape component is autocorrelated, so the iid value $\sigma_z^2/w$ understates $\var(\bar\varepsilon)$. For a stationary sequence with autocorrelation $\rho(k)$,
\begin{equation}\label{eq:app-tau}
  \var(\bar\varepsilon)
  = \frac{\sigma_z^2}{w} \sum_{|k| < w}\Big(1 - \frac{|k|}{w}\Big)\rho(k)
  \;\xrightarrow[w \gg 1]{}\; \frac{\sigma_z^2\, \tau}{w},
  \qquad
  \tau = \sum_{k=-\infty}^{\infty} \rho(k) = \frac{1+\rho_1}{1-\rho_1}
\end{equation}
for an AR(1) shape with lag-one autocorrelation $\rho_1$ ($\rho_1 = 0$ recovers the iid case); $\tau$ is the integrated autocorrelation time. Second, the ramp term $\sx$ is $\lambda$-coupled under this mapping (\Cref{rem:app-numeric}) and is evaluated numerically from the simulated covariate path rather than from the large-$w$ approximation, since window averaging enters through the correction factor $\bar\rho$. With these substitutions, plus $\sest$ measured out of sample from the first stage, the constrained-least-squares risks of Proposition~2$'$ become parameter-free predictions of the trained models' test MSE --- the adjudication standard used by the synthetic gate in \Cref{sec:synthetic}.

\subsection{Three crossover numbers, one anchor}\label{app:thresholds}

Three distinct values of $\lamstar$ appear in this paper and only one of them
anchors the pre-registered $\tau$; we set them side by side because conflating
them would make the threshold look either arbitrary or over-fitted.

\begin{center}
\fbox{\begin{minipage}{0.92\linewidth}
\textbf{Three thresholds, one anchor.} Three distinct crossover numbers
appear in this paper, and only one of them anchors $\tau$:
\begin{enumerate}
  \item $\lamstar_{\mathrm{M1}} = 0.923$ (at $h=24$; $0.664$ at $h=96$,
  $\lamstar<0$ at $h=336$) --- the \Cref{fig:dominance} \emph{baseline},
  computed with $\sx=0$ and $\sigma_\varepsilon=0$. Switching off
  covariate-driven level motion removes the channel that penalizes
  instance normalization within the horizon, so this is the most
  \inm-favorable stylization; it illustrates the geometry of the dominance
  map and is \emph{not} the anchor for $\tau$.
  \item $\lamstar_{\mathrm{M1}} \approx 0.27$--$0.28$ (at $h=24$) --- the
  same stylized model under the parameter mapping of the synthetic study
  (\Cref{sec:synthetic}), with covariate dynamics active. This is the
  restoration-rule \emph{upper bound} on the true crossover, and it is the
  number that anchors the pre-registered $\tau = 0.3$.
  \item $\lamstar_{\mathrm{OLS}} \approx 0.007$--$0.022$ --- the exact
  constrained-least-squares crossings (Proposition~2$'$), matched by
  trained RLinear models (empirical crossings $0.009$--$0.029$) to within
  $0.015$. This is where trained linear predictors actually cross.
\end{enumerate}
Ordering matters: $\tau$ sits at the stylized upper bound (ii), far above
the operative crossover (iii), which is what makes the decision rule
conservative toward instance normalization.
\end{minipage}}
\end{center}
